\subsection{The activation-tax construction}

The first attempt was to encode locality directly in the proximal geometry.
Let
\[
  h_\rho(x):=\frac12\norm{x}_2^2+\rho J(x).
\]
Given $y\ge0$ and a set $S\supseteq\supp(y)$, choose the particular
subgradient $g^S\in\partial J(y)$ by
\begin{equation}
  g_i^S:=
  \begin{cases}
    w_i,&i\in S,\\
    0,&i\notin S.
  \end{cases}
  \label{eq:selected-subgradient}
\end{equation}
The corresponding generalized Bregman divergence is
\[
 D_{h_\rho}^{g^S}(z,y)
 :=\frac12\norm{z-y}_2^2
   +\rho J(z)-\rho J(y)-\rho\ip{g^S}{z-y}.
\]

\begin{lemma}[Exact activation tax; \ProvedStatus]
\label{lem:activation-tax}
For $y,z\ge0$ with $\supp(y)\subseteq S$,
\begin{equation}
  D_{h_\rho}^{g^S}(z,y)
  =\frac12\norm{z-y}_2^2
   +\rho\sum_{i\notin S}w_i z_i.
  \label{eq:activation-tax}
\end{equation}
\end{lemma}

\begin{proof}
On the nonnegative orthant $J(z)=w^Tz$.  The linear terms cancel on $S$.
Outside $S$ we have $y_i=g_i^S=0$, leaving $\rho w_i z_i$.
\end{proof}

Thus the geometry is a Euclidean proximal term plus a linear reserve for every
coordinate not yet admitted to the face.

\subsection{What the matched subproblem proves}

Consider
\begin{equation}
 H_{y,S}(z)
 :=F_\rho(z)+\kappa D_{h_\rho}^{g^S}(z,y),
 \qquad \kappa>0.
 \label{eq:matched-subproblem}
\end{equation}
Up to constants,
\begin{equation}
 H_{y,S}(z)
 =\frac12z^T(Q+\kappa I)z
  -\ip{b+\kappa y+\kappa\rho g^S}{z}
  +(\alpha+\kappa)\rho J(z).
 \label{eq:matched-expansion}
\end{equation}

\begin{lemma}[Matched support budget; \ProvedStatus]
\label{lem:matched-budget}
Assume
\begin{equation}
  p(y)+\rho\vol(S)\le1.
  \label{eq:center-set-budget}
\end{equation}
Let $z^\star$ minimize~\eqref{eq:matched-subproblem}.  Then $z^\star\ge0$ and
\begin{equation}
 p(z^\star)+\rho\vol(\supp z^\star)\le1.
 \label{eq:matched-support-budget}
\end{equation}
\end{lemma}

\begin{proof}
The effective source in~\eqref{eq:matched-expansion} is nonnegative and
$Q+\kappa I$ is an $M$-matrix, hence $z^\star\ge0$.  Its stationarity equation
can be written
\[
 (Q+\kappa I)z^\star
 =b+\kappa y+\kappa\rho g^S
  -(\alpha+\kappa)\rho\,w\odot u,
\]
where $u_i=1$ when $z_i^\star>0$ and $u_i\in[0,1]$ otherwise.  Multiplication
by $w^T$ and~\eqref{eq:weighted-identities} gives
\[
 (\alpha+\kappa)
 \left(p(z^\star)+\rho\sum_i d_i u_i\right)
 =\alpha+\kappa\left(p(y)+\rho\vol(S)\right).
\]
Use~\eqref{eq:center-set-budget} and
$\sum_i d_i u_i\ge\vol(\supp z^\star)$.
\end{proof}

This is stronger than a post-hoc optimal support bound: every exact matched
inner minimizer is $1/\rho$-local.  A monotone coordinate solver started below
$z^\star$ remains inside the same order interval.

For the Catalyst shift, specialize the preceding parameter as
\begin{equation}
  \kappa=\kappa_{\mathrm A}:=1-2\alpha,
  \label{eq:kappa}
\end{equation}
the diagonal-to-weighted-eigenvalue ratio gives
\begin{equation}
 \chi_{\kappa_{\mathrm A}}
 :=\frac{2(\alpha+\kappa_{\mathrm A})}
         {1+\alpha+2\kappa_{\mathrm A}}
 =\frac23.
 \label{eq:inner-contraction}
\end{equation}
One complete sweep over currently positive KKT coordinates therefore contracts
their weighted mass by a constant factor.  If $\mathcal C_k^+$ is the positive
KKT mass at inner sweep $k$, then
\[
 \mathcal C_{k+1}^+\le\frac13\mathcal C_k^+.
\]
Moreover, for a monotone iterate $z\le z^\star$,
\[
 H_{y,S}(z)-H_{y,S}(z^\star)
 \le \pos{r(z)}^T(z^\star-z)
 \le \mathcal C^+(z).
\]
Consequently, an inner gap $\phi$ costs
\begin{equation}
 \Work_{\mathrm{inner}}(\phi)
 =O\!\left(
   \frac1\rho\log\frac{\mathcal C_0^+}{\phi}
 \right).
 \label{eq:matched-inner-work}
\end{equation}
The locality and inner-solution parts of the matched-Bregman proposal are
therefore valid.

\subsection{Exact face-change obstruction}

The failure occurs when accelerated interpolation activates a coordinate that
was zero at the center.  Take one coordinate with $y_i=0$, selected
subgradient $g_i^S=0$, and $z_i=a>0$.  By
\cref{lem:activation-tax},
\[
 D_{h_\rho}^{0}(ae_i,0)=\frac12a^2+\rho w_i a.
\]
For $\theta\in(0,1)$,
\begin{equation}
 D_{h_\rho}^{0}(\theta ae_i,0)
 =\theta^2D_{h_\rho}^{0}(ae_i,0)
  +\theta(1-\theta)\rho w_i a.
 \label{eq:face-change-defect}
\end{equation}

\begin{proposition}[No uniform exponent-two scaling; \RefutedStatus]
\label{prop:no-triangle-scaling}
There is no constant $G$ independent of $a$ and $\theta$ such that
\[
 D_{h_\rho}^{0}(\theta ae_i,0)
 \le G\theta^2D_{h_\rho}^{0}(ae_i,0)
\]
at every newly activated coordinate.
\end{proposition}

\begin{proof}
Divide~\eqref{eq:face-change-defect} by
$\theta^2D_{h_\rho}^{0}(ae_i,0)$ and let $a\downarrow0$.  The ratio tends to
$1/\theta$, so any admissible $G$ must satisfy $G\ge1/\theta$.
\end{proof}

The sharp face-change exponent is one, whereas Euclidean acceleration uses
exponent two.  This is precisely the distinction in the triangle-scaling
theory of \citet{hanzely2021accelerated}.  There is also a geometric warning:
on the positive orthant, the canonical derivative of
$h_\rho(x)=\frac12\norm{x}^2+\rho w^Tx$ produces a purely Euclidean Bregman
distance.  The activation tax exists only because the subgradient at zero is
chosen as a function of $S$.  It is a changing-face generalized divergence,
not one fixed smooth Bregman geometry.

\subsection{Why the defect does not telescope}

Let $S_t$ be the admitted face, $A_t=S_{t+1}\setminus S_t$, and define the
newly revealed optimal mass
\[
 \Gamma_t:=\sum_{i\in A_t}w_i x_{\bar\rho,i}^\star.
\]
Inserting~\eqref{eq:face-change-defect} into the usual accelerated
estimate-sequence template adds a term of the form
\begin{equation}
 \mathcal E_{t+1}
 \le q_{\mathrm{fc}}\mathcal E_t
     +C\theta\bar\rho\Gamma_t+e_t,
 \qquad
 q_{\mathrm{fc}}=1-c\sqrt\alpha,
 \qquad
 \theta=\Theta(\sqrt\alpha),
 \label{eq:defective-estimate-sequence}
\end{equation}
where $e_t$ is inner error.  Hence late activations carry essentially no
geometric discount.  The support budget gives only
$\sum_t\Gamma_t\le p(x_{\bar\rho}^\star)\le1$; it does not provide the much
smaller geometrically weighted mass needed to reach the handoff
gap~\eqref{eq:warm-gap}.

Adding the undiscovered mass
$R_t=\sum_{i\notin S_t}w_i x_{\bar\rho,i}^\star$ to the potential does not
repair the proof.  During a fixed-face accelerated step $R_t$ is unchanged,
so the combined recurrence contains a noncontracting term proportional to
$(1-q_{\mathrm{fc}})\bar\rho R_t$. The argument would have to assume that
unseen mass is
already small.

\subsection{Dual-envelope obstruction}

The conjugate route fails even in one dimension.  Let
\[
 h(x)=\frac12x^2+\rho|x|,
 \qquad
 F(x)=\frac\lambda2x^2-bx+\alpha\rho|x|,
\]
and consider
\[
 E_\kappa(u)
 :=\kappa h^\star(u)-(F+\kappa h)^\star(\kappa u).
\]
For $0\le u<\rho$ while the inner primal coordinate is active,
\[
 E_\kappa(u)
 =-\frac{[b+\kappa u-(\alpha+\kappa)\rho]^2}
         {2(\lambda+\kappa)},
 \qquad
 E_\kappa''(u)=-\frac{\kappa^2}{\lambda+\kappa}<0.
\]
Thus conjugation does not produce a globally convex envelope suitable for
projected Nesterov acceleration.  It becomes a regular quadratic only after
the active face has already been fixed.

\paragraph{No-go principle.}
Any generalized prox geometry that retains a positive linear reserve for
activation at zero has triangle-scaling exponent one at a new coordinate.
Choosing the canonical derivative removes the defect and the activation tax
simultaneously.  Changing constants or replacing $J$ by another locally
linear sparsity potential cannot resolve this conflict.
